% ===========================================================================
% 05 - Live-Demo, FR/QR-Nachweise und Fazit
% ===========================================================================

\begin{frame}[fragile]{Live-Demo -- Circuit Breaker und Retry}
  \frbadge{FR-10/11}
  \begin{columns}[T,onlytextwidth]
    \begin{column}{0.5\textwidth}
      \textbf{Ablauf der Demo}
      \begin{enumerate}
        \item System starten:
          \begin{lstlisting}[style=bash]
docker compose up --build
          \end{lstlisting}
        \item Thermal stoppen:
          \begin{lstlisting}[style=bash]
docker compose stop thermal-analysis-service
          \end{lstlisting}
        \item Analyse starten $\rightarrow$ \code{THERMAL = FAILED}.
        \item Thermal starten:
          \begin{lstlisting}[style=bash]
docker compose start thermal-analysis-service
          \end{lstlisting}
        \item Retry ausführen $\rightarrow$ \code{overallResult = OK}.
      \end{enumerate}
    \end{column}
    \begin{column}{0.46\textwidth}
      \textbf{Retry-Endpoint}
      \begin{lstlisting}[style=bash]
POST /api/analyses/{analysisId}/algorithms/THERMAL/retry
      \end{lstlisting}

      \vspace{0.1cm}
      \textbf{Was die Demo zeigt}
      \begin{itemize}
        \item \textbf{TA-04 / ADR-04}: Circuit Breaker fängt den Ausfall ab.
        \item \textbf{FR-11}: System bleibt erreichbar, Schritt wird \code{FAILED}.
        \item \textbf{FR-10}: Nur der fehlgeschlagene Schritt wird wiederholt.
        \item \textbf{ADR-03}: Choreographie läuft ab Thermal weiter.
      \end{itemize}
    \end{column}
  \end{columns}
\end{frame}

\begin{frame}[fragile]{FR-07 -- Gesamtergebnis im Code}
  \frbadge{FR-07}
  \begin{columns}[T,onlytextwidth]
    \begin{column}{0.54\textwidth}
      \begin{lstlisting}[style=java]
public void recalculateOverallResult() {
    boolean failed = executions.stream()
        .anyMatch(e -> e.getStatus() == AnalysisStatus.FAILED
                    || e.getResult() == AnalysisResult.FAILED);

    if (failed) {
        overallResult = AnalysisResult.FAILED;
        return;
    }

    boolean unfinished = executions.stream()
        .anyMatch(e -> e.getStatus() == AnalysisStatus.PENDING
                    || e.getStatus() == AnalysisStatus.RUNNING);

    overallResult = unfinished ? null : AnalysisResult.OK;
}
      \end{lstlisting}
      \srcfile{analysis-management-service/.../domain/AnalysisRun.java}
    \end{column}
    \begin{column}{0.42\textwidth}
      \textbf{Regel (architecture.md \S8)}
      \begin{itemize}
        \item Ein Fehler $\rightarrow$ \code{FAILED}.
        \item Schritte offen $\rightarrow$ \code{null}.
        \item Alle vier fertig $\rightarrow$ \code{OK}.
      \end{itemize}
      \vspace{0.2cm}
      \begin{block}{Test}
        \small \code{AnalysisRunTest} prüft alle drei Fälle (QR-07).
      \end{block}
    \end{column}
  \end{columns}
\end{frame}

\begin{frame}[fragile]{QR-08 -- Responsiveness durch \code{@Async}}
  \qrbadge{QR-08}
  \begin{columns}[T,onlytextwidth]
    \begin{column}{0.5\textwidth}
      \begin{lstlisting}[style=java]
@PostMapping
public ResponseEntity<Void> start(
        @RequestBody AnalysisCommand command) {
    worker.execute(command);
    return ResponseEntity.accepted().build();
}
      \end{lstlisting}
      \srcfile{.../api/AnalysisController.java}

      \vspace{0.1cm}
      \begin{lstlisting}[style=java]
@Async
public void execute(AnalysisCommand command) {
    // laeuft im Hintergrund-Thread
}
      \end{lstlisting}
      \srcfile{.../application/AnalysisWorker.java}
    \end{column}
    \begin{column}{0.46\textwidth}
      \textbf{Erläuterung}
      \begin{itemize}
        \item Der Controller antwortet sofort mit \code{202 Accepted}.
        \item \code{worker.execute} läuft dank \code{@Async} im Hintergrund.
        \item Der Client blockiert nicht bis zum Ende der Analyse.
        \item Status ist währenddessen weiter abfragbar (QR-02).
      \end{itemize}
      \begin{block}{Zusammenspiel}
        \small \code{@EnableAsync} (TA-01) + \code{@Async} + \code{202 Accepted}.
      \end{block}
    \end{column}
  \end{columns}
\end{frame}

\begin{frame}{Fazit -- Anforderungen nachweisbar umgesetzt}
  \begin{columns}[T,onlytextwidth]
    \begin{column}{0.5\textwidth}
      \textbf{Technische Anforderungen}
      \begin{itemize}
        \item[\textcolor{okgreen}{\checkmark}] TA-01 Spring Boot
        \item[\textcolor{okgreen}{\checkmark}] TA-02 Docker
        \item[\textcolor{okgreen}{\checkmark}] TA-03 Docker Compose
        \item[\textcolor{okgreen}{\checkmark}] TA-04 Resilience4j
        \item[\textcolor{codegray}{--}] TA-05 Kafka (COULD, bewusst offen)
      \end{itemize}

      \vspace{0.3cm}
      \textbf{Architekturentscheidungen}
      \begin{itemize}
        \item[\textcolor{okgreen}{\checkmark}] ADR-01 bis ADR-06 im Code belegt
      \end{itemize}
    \end{column}
    \begin{column}{0.46\textwidth}
      \begin{block}{Kernbotschaft}
        \small Jede zentrale Anforderung ist über eine konkrete Datei und ein
        Symbol im Code nachweisbar -- nicht nur dokumentiert, sondern
        implementiert und testbar.
      \end{block}

      \vspace{0.2cm}
      \begin{block}{Referenzen}
        \small
        \code{docs/code-traceability.md}\\
        \code{docs/arc42.md}\\
        \code{docs/requirements.md}\\
        \code{scripts/e2e.sh}
      \end{block}
    \end{column}
  \end{columns}
\end{frame}

\begin{frame}[plain]
  \begin{center}
    \hbrslogo
    \vspace{1.5cm}

    {\Huge\bfseries\color{darkblue} Vielen Dank!}\par
    \vspace{0.6cm}
    {\large Fragen und Diskussion}\par

    \vspace{1.5cm}
    {\large Theodoros Bampis \& Yunus Emre Sirin}\par
    {\large SEKA -- Sommersemester 2026}\par
  \end{center}
\end{frame}
